\documentclass[a4paper,12pt]{article}
\usepackage[english]{babel}
\usepackage{hyperref}
\hypersetup{hidelinks}
\usepackage{amsmath}
\usepackage{amssymb}
\usepackage{physics}
\usepackage{bm}
\usepackage[varvw]{newtxmath}
\usepackage{booktabs}
\usepackage{array}

% page geometry — identical to thesis
\setlength{\oddsidemargin}{8mm}
\setlength{\textwidth}{145mm}
\setlength{\topmargin}{2mm}
\setlength{\textheight}{225mm}
\setlength{\headheight}{0mm}
\setlength{\parskip}{0.5mm}
\setlength{\parindent}{9mm}
\linespread{1.12}
\renewcommand{\floatpagefraction}{.93}
\renewcommand{\textfraction}{0.02}

% ── convenience macros ───────────────────────────────────────────────────────
\newcommand{\dt}{\Delta t}
\newcommand{\dx}{\Delta x}
\newcommand{\Hhat}{\hat{H}}

% ── title block ─────────────────────────────────────────────────────────────
\title{{\bf \Large
    Numerical Simulation of 1D Quantum Tunnelling\\[2mm]
    via the Time-Dependent Schr\"{o}dinger Equation}}
\author{Ludovico Fecia di Cossato\\[1mm]}
\date{March 2026}

% ─────────────────────────────────────────────────────────────────────────────
\begin{document}

\maketitle
\thispagestyle{empty}

\begin{center}
  \small\textit{Atomic units throughout: $\hbar = 1$, $m = 1$.}
\end{center}

\vskip 50mm
\noindent\rule{\textwidth}{0.4pt}

% ── Abstract ─────────────────────────────────────────────────────────────────
\begin{abstract}
This document describes the design and implementation of a one-dimensional
solver for the Time-Dependent Schr\"{o}dinger Equation (TDSE) based on the
Crank--Nicolson finite-difference scheme.
The physical objective is to simulate the time evolution of a Gaussian wave
packet interacting with a rectangular potential barrier, resolving the quantum
tunnelling phenomenon and extracting the transmission and reflection
coefficients as functions of the barrier parameters.
The solver is implemented in C++17 with LAPACKE for the tridiagonal complex
linear system; post-processing is handled by Python scripts.
\end{abstract}

\noindent\rule{\textwidth}{0.4pt}
\vskip 4mm

\newpage

% ─────────────────────────────────────────────────────────────────────────────
\section{Physical Background}

\subsection{The wave function}

The state of a non-relativistic quantum particle in one spatial dimension is
completely described by a complex-valued wave function $\psi(x,t)$.
Its squared modulus $|\psi(x,t)|^2$ is the probability density of finding the
particle at position $x$ at time $t$.
Physical consistency requires the normalisation condition to hold at all times:
\begin{equation}
  \int_{-\infty}^{+\infty} |\psi(x,t)|^2 \, dx = 1.
\end{equation}
Any admissible time-evolution operator must therefore be \emph{unitary},
preserving the norm of $\psi$ exactly.
This provides the primary numerical verification criterion: if the norm drifts,
the solver contains an error.

\subsection{The Time-Dependent Schr\"{o}dinger Equation}

In atomic units ($\hbar = 1$, $m = 1$), the TDSE in one spatial dimension reads:
\begin{equation}
  i\,\frac{\partial\psi}{\partial t}
  = \hat{H}\psi
  = \left[-\frac{1}{2}\frac{\partial^2}{\partial x^2} + V(x)\right]\psi.
\end{equation}
The left-hand side governs time evolution; the right-hand side is the
Hamiltonian operator $\hat{H}$, composed of the kinetic energy operator
$-\tfrac{1}{2}\partial^2/\partial x^2$ and the potential energy $V(x)$ acting
multiplicatively.
The formal solution is $\psi(x,t) = e^{-i\hat{H}t}\,\psi(x,0)$, where the
time-evolution operator cannot be computed analytically for a general $V(x)$,
motivating the numerical approach.

\subsection{The Gaussian wave packet}

The most natural initial condition for a scattering problem is a
\emph{Gaussian wave packet} — a spatially localised state with a well-defined
mean momentum.
In atomic units:
\begin{equation}
  \psi(x,0)
  = \bigl(2\pi\sigma^2\bigr)^{-1/4}
    \exp\!\left[-\frac{(x-x_0)^2}{4\sigma^2}\right]
    \exp(ik_0 x).
\end{equation}
The three parameters have clear physical interpretations.
$x_0$ is the initial position of the packet centre;
$\sigma$ controls its spatial width (and inversely its momentum spread, via the
Heisenberg uncertainty relation $\Delta x\,\Delta p \ge \tfrac{1}{2}$);
$k_0$ is the central wave vector, setting the mean momentum $\langle p\rangle = k_0$
and therefore the direction and speed of propagation.
The prefactor ensures unit normalisation.

\subsection{The rectangular potential barrier}

The potential studied here is a rectangular barrier of height $V_0$, left edge
at $x_\mathrm{bar}$, and width $w_\mathrm{bar}$:
\begin{equation}
  V(x) =
  \begin{cases}
    V_0 & x_\mathrm{bar} \le x \le x_\mathrm{bar} + w_\mathrm{bar},\\
    0   & \text{otherwise.}
  \end{cases}
\end{equation}
Classically, a particle with kinetic energy $E < V_0$ is always reflected.
Quantum mechanically, the wave function has a non-zero amplitude inside and
beyond the barrier, leading to a finite transmission probability $T > 0$.
This is \emph{quantum tunnelling} — one of the most striking departures from
classical mechanics and the central phenomenon this project aims to simulate.
The transmission probability decreases exponentially with the barrier width
$w_\mathrm{bar}$, making it a sensitive probe of the solver's accuracy.

\subsection{Observables}

Once the time evolution of $\psi$ is known, the physically relevant quantities
are extracted as expectation values or integrated quantities:

\vskip 2mm
\begin{center}
\begin{tabular}{lll}
  \toprule
  Observable & Definition & Physical meaning \\
  \midrule
  Norm        & $\int|\psi|^2\,dx$                  & Must equal 1 at all times \\
  $\langle x\rangle(t)$ & $\int x|\psi|^2\,dx$     & Mean position of the particle \\
  $T$         & $\int_{x > x_\mathrm{bar}+w_\mathrm{bar}}|\psi|^2\,dx$ & Transmission probability \\
  $R$         & $\int_{x < x_\mathrm{bar}}|\psi|^2\,dx$ & Reflection probability \\
  \bottomrule
\end{tabular}
\end{center}
\vskip 2mm

\noindent
A well-functioning solver must satisfy $T + R \approx 1$ at all times —
probability conservation expressed in terms of scattering coefficients.

% ─────────────────────────────────────────────────────────────────────────────
\section{Space-Time Discretisation}

The TDSE is a partial differential equation in two continuous variables $x$ and
$t$.
To solve it numerically, both must be replaced by a finite grid.

\subsection{Spatial grid}

The physical domain $[0, L]$ is divided into $N$ equal subintervals of width
$\dx = L/N$.
The grid points are $x_j = j\dx$ for $j = 0, 1, \ldots, N$.
At any fixed time $t$, the wave function becomes a vector of $N+1$ complex
numbers:
\begin{equation}
  \psi_j(t) \;\equiv\; \psi(x_j,\,t), \qquad j = 0,1,\ldots,N.
\end{equation}
The boundary conditions $\psi(0,t) = \psi(L,t) = 0$ model an infinite potential
well — the particle is confined to the domain.
This fixes the two endpoints and leaves $N-1$ internal degrees of freedom to
evolve.

\subsection{Temporal grid}

Time is discretised with a uniform step $\dt$.
The simulation runs from $t=0$ to $t = T_\mathrm{tot}$ in
$M = T_\mathrm{tot}/\dt$ steps.
The state at step $n$ is written $\psi_j^n \equiv \psi(x_j, t_n)$.
The central question becomes: given $\psi^n$ (the state now), how is
$\psi^{n+1}$ (the state one step later) computed?

\subsection{Finite-difference approximation of derivatives}

The time derivative is approximated by a forward difference:
\begin{equation}
  \frac{\partial\psi}{\partial t}\bigg|_{x_j,\,t_n}
  \approx \frac{\psi_j^{n+1} - \psi_j^n}{\dt}.
\end{equation}
The second spatial derivative is approximated by the standard centred
three-point formula:
\begin{equation}
  \frac{\partial^2\psi}{\partial x^2}\bigg|_{x_j}
  \approx \frac{\psi_{j-1} - 2\psi_j + \psi_{j+1}}{\dx^2},
\end{equation}
which has truncation error $\mathcal{O}(\dx^2)$, as verified by Taylor
expansion.
The key remaining choice — at which time level to evaluate the spatial
derivative — determines the stability and unitarity of the scheme, and is the
subject of the next section.

% ─────────────────────────────────────────────────────────────────────────────
\section{The Crank--Nicolson Scheme}

\subsection{Motivation and construction}

A purely explicit scheme evaluates the spatial derivative at time $t_n$;
a purely implicit scheme evaluates it at $t_{n+1}$.
Neither alone preserves the norm of $\psi$.
The \emph{Crank--Nicolson} scheme takes the average of both, effectively
evaluating the Hamiltonian at the midpoint $t_{n+1/2}$.
Substituting into the discretised TDSE:
\begin{equation}
  i\,\frac{\psi^{n+1} - \psi^n}{\dt}
  = \frac{1}{2}\hat{H}_\mathrm{disc}\bigl(\psi^{n+1} + \psi^n\bigr).
\end{equation}
Rearranging, one obtains a linear system to be solved at each time step:
\begin{equation}
  A\,\psi^{n+1} = B\,\psi^n,
\end{equation}
where the matrices $A$ and $B$ are:
\begin{equation}
  A = I + \frac{i\dt}{2}\hat{H}_\mathrm{disc}, \qquad
  B = I - \frac{i\dt}{2}\hat{H}_\mathrm{disc}.
\end{equation}
Note that $B = A^\dagger$ (the conjugate transpose of $A$).
This immediately implies that the scheme is \emph{unitary}:
$\|\psi^{n+1}\| = \|\psi^n\|$ exactly, independent of the choice of $\dt$
and $\dx$.

\subsection{Tridiagonal structure and solution}

Because the finite-difference kinetic operator only couples each point to its
two neighbours, both $A$ and $B$ are \emph{tridiagonal} matrices of size
$(N-1)\times(N-1)$.
Setting $r = i\dt/(4\dx^2)$, the diagonals are:
\begin{align}
  A:\quad & \text{main diagonal} = 1 + 2r + \tfrac{i\dt}{2}V_j, \qquad
            \text{off-diagonals} = -r, \notag\\
  B:\quad & \text{main diagonal} = 1 - 2r - \tfrac{i\dt}{2}V_j, \qquad
            \text{off-diagonals} = +r. \notag
\end{align}
At each time step, the right-hand side vector
$\mathbf{b}^n = B\,\psi^n$ is computed by a matrix-vector multiplication,
and the tridiagonal system $A\,\psi^{n+1} = \mathbf{b}^n$ is solved using the
LAPACK routine \texttt{LAPACKE\_zgtsv}, which operates on
\texttt{std::complex<double>} arrays and exploits the tridiagonal structure
for $\mathcal{O}(N)$ cost per time step.

% ─────────────────────────────────────────────────────────────────────────────
\section{Project Structure and Parameters}

\subsection{Implementation blocks}

The solver is organised in five logical blocks, each building on the previous:

\vskip 2mm
\begin{center}
\begin{tabular}{ll}
  \toprule
  Block & Content \\
  \midrule
  0 — Physics        & TDSE, Hamiltonian, observables \\
  1 — Discretisation & Spatial and temporal grid, $\psi$ as a complex vector \\
  2 — Crank--Nicolson & Implicit/explicit averaging, matrices $A$ and $B$ \\
  3 — Linear system  & Tridiagonal structure, \texttt{LAPACKE\_zgtsv} solver \\
  4 — Initial cond.  & Gaussian wave packet, boundary conditions $\psi = 0$ \\
  5 — Observables    & Norm, $\langle x\rangle$, $T$, $R$ — verification and output \\
  \bottomrule
\end{tabular}
\end{center}
\vskip 2mm

\subsection{Source layout}

The C++17 codebase follows a clean separation between data, logic, and
execution:
\begin{verbatim}
TDSE_CN/
|-- input/input.txt        # physical and numerical parameters
|-- src/params.hpp/.cpp    # Params struct, file I/O, printing
|-- src/physics.hpp/.cpp   # grid, wave packet, matrices, time step
|-- src/tdse.cpp           # main: orchestration
|-- output/                # generated data (gitignored)
|-- analysis/              # Python plotting scripts
|-- docs/                  # figures and reports
+-- makefile
\end{verbatim}

\subsection{Input parameters}

All physical and numerical parameters are read from a plain-text file at
runtime, with no hardcoded values in the source.
Parameters are listed below with their units in the atomic system
($\hbar = 1$, $m = 1$):

\vskip 2mm
\begin{center}
\begin{tabular}{llll}
  \toprule
  Parameter & Symbol & Unit & Description \\
  \midrule
  \texttt{L}     & $L$               & bohr          & Total domain length \\
  \texttt{dx}    & $\dx$             & bohr          & Spatial step \\
  \texttt{dt}    & $\dt$             & $\hbar/E_h$   & Time step \\
  \texttt{T\_tot}& $T_\mathrm{tot}$  & $\hbar/E_h$   & Total simulation time \\
  \texttt{x0}    & $x_0$             & bohr          & Initial packet centre \\
  \texttt{sigma} & $\sigma$          & bohr          & Packet width \\
  \texttt{k0}    & $k_0$             & bohr$^{-1}$   & Initial wave vector \\
  \texttt{x\_bar}& $x_\mathrm{bar}$  & bohr          & Left edge of barrier \\
  \texttt{w\_bar}& $w_\mathrm{bar}$  & bohr          & Barrier width \\
  \texttt{V0}    & $V_0$             & $E_h$         & Barrier height \\
  \bottomrule
\end{tabular}
\end{center}
\vskip 2mm

\noindent
The derived quantities $N = \mathrm{round}(L/\dx)$ and
$M = \mathrm{round}(T_\mathrm{tot}/\dt)$ are computed automatically from the
input using \texttt{std::round} to avoid off-by-one errors from floating-point
division.

\subsection{Parameter selection guidelines}

The input values are not arbitrary — a physically meaningful simulation
requires the following hierarchy.
The packet width $\sigma$ must be large enough to constitute a well-defined
localised state but small enough to fit away from the domain boundaries at
$t=0$.
The initial position $x_0$ must be far enough from the barrier to allow the
packet to form before impact.
The kinetic energy $E \sim k_0^2/2$ should be comparable to $V_0$ to produce
interesting tunnelling behaviour: if $E \gg V_0$ the barrier is transparent;
if $E \ll V_0$ essentially nothing tunnels.
Although Crank--Nicolson is unconditionally stable, $\dt$ must be chosen small
enough for accuracy.

\end{document}